\section{ER Diagrams}
This section taught the basics of \newterm{entity-relationship diagrams} (\index{ER Diagrams|see{Entity-Relationship Diagrams}}) and how to use them to model a database.

Not only did we learn how to diagram, but also Lab 2 emphasized coding our own, which was done in the programming language \newterm{mermaid}, which fits directly into LaTeX. 

Below is a simple example of 3 different types of relationships, one-to-one, one-to-many, and many-to-many.
% requires the mermaid package to be included in the preamble of main.tex
% and to be compiled with lualatex or xelatex, not pdflatex.
% needs shell escape
\begin{mermaid}
erDiagram
    %% One-to-one
    PERSON ||--|| PASSPORT : "has (one to one)"

    %% One-to-many
    DEPARTMENT ||--o{ EMPLOYEE : "employs (one to many)"

    %% Many-to-many
    STUDENT {
        int student_id PK
        string name
        string email
    }
    COURSE {
        int course_id PK
        string title
    }
    STUDENT }o--o{ COURSE : "enrolls_in (zero or many to zero or many)"
\end{mermaid}

Note that the symbols next to a node $A$ represent how many of $A$ can be in $B$. For example, a department has many or zero employees, but an employee must have exactly one department.